\documentclass[a4paper,9pt,twocolumn,twoside,]{pinp}

%% Some pieces required from the pandoc template
\providecommand{\tightlist}{%
  \setlength{\itemsep}{0pt}\setlength{\parskip}{0pt}}

% Use the lineno option to display guide line numbers if required.
% Note that the use of elements such as single-column equations
% may affect the guide line number alignment.

\usepackage[T1]{fontenc}
\usepackage[utf8]{inputenc}

% pinp change: the geometry package layout settings need to be set here, not in pinp.cls
\geometry{layoutsize={0.95588\paperwidth,0.98864\paperheight},%
  layouthoffset=0.02206\paperwidth, layoutvoffset=0.00568\paperheight}

\definecolor{pinpblue}{HTML}{185FAF}  % imagecolorpicker on blue for new R logo
\definecolor{pnasbluetext}{RGB}{101,0,0} %



\title{Women in Parliament -- Tidyverse Edition}

\author[]{Saghir Bashir}


\setcounter{secnumdepth}{3}

% Please give the surname of the lead author for the running footer
\leadauthor{\href{https://creativecommons.org/licenses/by-sa/4.0/}{CC BY SA}
\href{https://ilustat.com/}{ilustat} \(\bullet\)
\href{mailto:info@ilustat.com}{\nolinkurl{info@ilustat.com}}}

% Keywords are not mandatory, but authors are strongly encouraged to provide them. If provided, please include two to five keywords, separated by the pipe symbol, e.g:
 \keywords{  Women in Parliament |  World Bank Indicator |  Tidyverse |  dplyr |  tidyr |  ggplot2  }  

\begin{abstract}
We will use the World Bank's indicator data for ``Women in Parliament''
as a case study when working with the tidyverse suite of R packages. We
will guide you through the geographical and time trends for the
percentage of women in national parliaments. We will start by learning
about and understanding the raw data, which we will then process
(``wrangle'') in preparation for some exploratory analysis.
\end{abstract}

\dates{This version was compiled on \today} 

% initially we use doi so keep for backwards compatibility
% new name is doi_footer
\doifooter{Learn more at \url{https://ilustat.com/resources/}}

\pinpfootercontents{Women in Parliament -- Tidyverse}

\begin{document}

% Optional adjustment to line up main text (after abstract) of first page with line numbers, when using both lineno and twocolumn options.
% You should only change this length when you've finalised the article contents.
\verticaladjustment{-2pt}

\maketitle
\thispagestyle{firststyle}
\ifthenelse{\boolean{shortarticle}}{\ifthenelse{\boolean{singlecolumn}}{\abscontentformatted}{\abscontent}}{}

% If your first paragraph (i.e. with the \dropcap) contains a list environment (quote, quotation, theorem, definition, enumerate, itemize...), the line after the list may have some extra indentation. If this is the case, add \parshape=0 to the end of the list environment.


\begin{tikzpicture}[remember picture, overlay]
    \node [shift={(-2cm,-2cm)}]  at (current page.north east)
        {\includegraphics[width=2.5cm]{../images/Women_in_Parliament_hex.png}};
\end{tikzpicture}

\hypertarget{preface}{%
\section{Preface}\label{preface}}

We present a real-life case study for some of the
\texttt{tidyverse}\footnote{For more information on the
  \texttt{tidyverse} see \url{https://www.tidyverse.org/}.} package
using the World Bank's ``Women in Parliament'' indicator data. To get
the most out of this case-study guide, repeat the examples and do the
exercises whilst reading it.

\hypertarget{guide-materials}{%
\subsection{Guide materials}\label{guide-materials}}

You can download materials for this guide from this link:

\begin{itemize}
\tightlist
\item
  \url{https://ilustat.com/shared/WiP-tidyverse.zip}
\end{itemize}

Unzip the file, which contains the data, this guide and an R script
exercise file. We advise you to work with ``\texttt{WiP-Exercise.R}''
file to follow the examples and do the exercises. If you are using
RStudio, you can double click on ``\texttt{WiP-tv.Rproj}'' to get
started.

\hypertarget{objectives}{%
\section{Objectives}\label{objectives}}

\emph{Explore the geographical and time trends for the percentage of
women\footnote{The objective could be termed neutrally as ``gender
  trends'' but we will keep it per the World Bank data.} in national
parliaments.}

\hypertarget{understanding-the-data}{%
\section{Understanding the Data}\label{understanding-the-data}}

\hypertarget{the-world-bank-data}{%
\subsection{The World Bank Data}\label{the-world-bank-data}}

The raw data for \emph{``Proportion of seats held by women in national
parliaments''} includes the percentage of women in parliament
(\emph{``single or lower parliamentary chambers only''}) by country
(region) and year. It can be downloaded from:\footnote{The
  \texttt{wbstats} R package
  (\url{https://cran.r-project.org/web/packages/wbstats/}) gives access
  to a ``tidier'' version of the World Bank indicator data.}

\begin{itemize}
\tightlist
\item
  \url{https://data.worldbank.org/indicator/SG.GEN.PARL.ZS}
\end{itemize}

As part of its ``open data'' mission the World Bank offers \emph{``free
and open access to global development data''} kindly licensed under the
``Creative Commons Attribution 4.0 (CC-BY 4.0)''.\footnote{\url{https://datacatalog.worldbank.org/public-licenses\#cc-by}.}

\hypertarget{source-data}{%
\subsection{Source Data}\label{source-data}}

The data originates from the ``Inter-Parliamentary Union''
(IPU)\footnote{Inter-Parliamentary Union: \url{https://www.ipu.org/}.}
which provides an \emph{\emph{``Archive of statistical data on the
percentage of women in national parliaments''}} going back to 1997 on a
monthly basis:

\begin{itemize}
\tightlist
\item
  \url{http://archive.ipu.org/wmn-e/classif-arc.htm}
\end{itemize}

The World Bank data is for ``single or lower parliamentary chambers
only'', while the IPU also presents data for ``Upper Houses or
Senates''. Moreover, the IPU provides the actual numbers used to
calculate the percentages (which the World Bank does not).

\hypertarget{data-limitations}{%
\subsection{Data limitations}\label{data-limitations}}

Take caution when interpreting these data, as parliamentary systems vary
from country to country, and in some cases over time. Some of the issues
to consider include:

\begin{itemize}
\tightlist
\item
  Who has, and who does not have, the right to become a Member of
  Parliament (MP)?
\item
  How does someone become an MP? Through democratic elections? How is
  ``democratic election'' defined?
\item
  What is the real power of MPs and their parliament? Can MPs make a
  difference?
\end{itemize}

\hypertarget{data-definitions-assumptions}{%
\subsection{Data definitions \&
assumptions}\label{data-definitions-assumptions}}

\hypertarget{women}{%
\subsubsection{``Women''}\label{women}}

The definition for ``women'' is not given, so we will assume that it
refers to a binary classification for gender (sex).

\hypertarget{country-region}{%
\subsubsection{``Country (Region)''}\label{country-region}}

The definition of countries and regions can change over time.
(e.g.~formation of new countries after conflicts, new member states
joining a pre-existing collective). How are these changes reflected in
the data? How do they affect the interpretation?

\hypertarget{pro-tip}{%
\subsection{Pro tip}\label{pro-tip}}

Understand the limitations of your data before anybody else points them
out to you.

\hypertarget{about-the-data-file}{%
\section{About the data file}\label{about-the-data-file}}

The data is stored in a file called:

\begin{itemize}
\tightlist
\item
  \texttt{API\_SG.GEN.PARL.ZS\_DS2\_en\_csv\_v2\_382411.csv}
\end{itemize}

To simplify things we have copied it to \texttt{WB-WiP.csv} (which also
allows us to maintain the original file in case something goes wrong).

\hypertarget{pro-tip-1}{%
\subsection{Pro tip}\label{pro-tip-1}}

Always keep a backup copy of the data. Alternatively, set the data
file(s) to ``read-only'' to protect it from being overwritten or
modified.

\hypertarget{exercise}{%
\subsection{Exercise}\label{exercise}}

It is important to look at and understand the contents of the file
before you start using it. Using a text editor or a spreadsheet
software, open the \texttt{WB-WiP.csv} file (in the \texttt{data}
directory). What do you observe in the contents of this file?

\hypertarget{content-and-structure}{%
\subsection{Content and Structure}\label{content-and-structure}}

The first four lines of WB-WiP.csv can be ignored, since they contain
two lines of meta-information and two blank lines, as follows:

\begin{verbatim}
 1  "Data Source","World Development Indicators",
 2  
 3  "Last Updated Date","2019-10-16",
 4  
\end{verbatim}

The fifth line contains the column (variable) names and the body of data
starts in the sixth line. It is important to note that there was no
collection of data for a majority of the years, which means that it is
``missing''.

\hypertarget{loading-tidyverse-packages}{%
\section{\texorpdfstring{Loading \texttt{tidyverse}
packages}{Loading tidyverse packages}}\label{loading-tidyverse-packages}}

We will load the \texttt{tidyverse} packages we plan to use individually
(messages have been suppressed).\footnote{We could have used
  \texttt{library(tidyverse)} but we prefer to load packages
  individually and only those that we will use.}

\begin{Shaded}
\begin{Highlighting}[]
\KeywordTok{library}\NormalTok{(here)}
\KeywordTok{library}\NormalTok{(readr)}
\KeywordTok{library}\NormalTok{(dplyr)}
\KeywordTok{library}\NormalTok{(tidyr)}
\KeywordTok{library}\NormalTok{(ggplot2)}
\KeywordTok{library}\NormalTok{(gghighlight)}
\end{Highlighting}
\end{Shaded}

\hypertarget{importing-the-data}{%
\section{Importing the data}\label{importing-the-data}}

Based on our findings above, we can ``skip'' the first four lines and
treat the fifth line as column (variable) names.

\begin{Shaded}
\begin{Highlighting}[]
\NormalTok{wip <-}\StringTok{ }\KeywordTok{read_csv}\NormalTok{(}\KeywordTok{here}\NormalTok{(}\StringTok{"data"}\NormalTok{, }\StringTok{"WB-WiP.csv"}\NormalTok{), }
             \DataTypeTok{skip =} \DecValTok{4}\NormalTok{)}
\end{Highlighting}
\end{Shaded}

\begin{ShadedResult}
\begin{verbatim}
#  Warning: Missing column names filled in:
#  'X64' [64]
\end{verbatim}
\end{ShadedResult}

\hypertarget{messages}{%
\subsubsection{Messages}\label{messages}}

We have suppressed some of the messages but left the one about variable
\texttt{X64} (we will come back to it below).

\hypertarget{exercise-1}{%
\subsection{Exercise}\label{exercise-1}}

Check what you have read by typing ``\texttt{wip}'' in the console
window. What do you observe? Type ``\texttt{class(wip)}'' and
``\texttt{glimpse(wip)}'' to confirm that ``\texttt{wip}'' is of class
``\texttt{tbl\_df}''.

\hypertarget{fix-column-names}{%
\subsection{``Fix'' column names}\label{fix-column-names}}

Some of the column names contain spaces while others are numeric:

\begin{Shaded}
\begin{Highlighting}[]
\KeywordTok{head}\NormalTok{(}\KeywordTok{names}\NormalTok{(wip))}
\CommentTok{#  [1] "Country Name"   "Country Code"  }
\CommentTok{#  [3] "Indicator Name" "Indicator Code"}
\CommentTok{#  [5] "1960"           "1961"}
\KeywordTok{tail}\NormalTok{(}\KeywordTok{names}\NormalTok{(wip))}
\CommentTok{#  [1] "2014" "2015" "2016" "2017" "2018" "X64"}
\end{Highlighting}
\end{Shaded}

By using the \texttt{make.names()} function we don't need to use back
ticks (\texttt{\textasciigrave{}}) around the column names
(e.g.~\texttt{\textasciigrave{}col\ name\textasciigrave{}}).

\begin{Shaded}
\begin{Highlighting}[]
\KeywordTok{names}\NormalTok{(wip) <-}\StringTok{ }\KeywordTok{make.names}\NormalTok{(}\KeywordTok{names}\NormalTok{(wip))}
\KeywordTok{head}\NormalTok{(}\KeywordTok{names}\NormalTok{(wip))}
\CommentTok{#  [1] "Country.Name"   "Country.Code"  }
\CommentTok{#  [3] "Indicator.Name" "Indicator.Code"}
\CommentTok{#  [5] "X1960"          "X1961"}
\KeywordTok{tail}\NormalTok{(}\KeywordTok{names}\NormalTok{(wip))}
\CommentTok{#  [1] "X2014" "X2015" "X2016" "X2017" "X2018"}
\CommentTok{#  [6] "X64"}
\end{Highlighting}
\end{Shaded}

\hypertarget{data-wrangling-aims}{%
\section{Data Wrangling Aims}\label{data-wrangling-aims}}

We can simplify the production of summaries and plots by restructuring
the current \texttt{wip} dataset (which has 64 columns) to the following
format:

\begin{verbatim}
   Country       Year   pctWiP
   Country AAA   1997     ##.#
   Country AAA   1998     ##.#
   Country AAA   1999     ##.#
   ...
\end{verbatim}

pctWiP refers to the percentage of women in parliament.

\hypertarget{key-information-retained}{%
\subsection{Key information retained}\label{key-information-retained}}

These three columns will contain the same information as the
\texttt{wip} dataset but in a more usable format. We will also add a
variable for the ratio of male to female MPs.

\hypertarget{superfluous-columns}{%
\subsection{Superfluous columns}\label{superfluous-columns}}

We will start by removing columns \texttt{X64}, \texttt{Indicator.Name}
and \texttt{Indicator.Code}. There are years without any data but they
will be removed automatically later (when restructuring from ``wide'' to
``long'' format).

Column \texttt{X64} is created automatically due to an extra comma at
the end of the column names (fifth) line of \texttt{WB-WiP.csv}:

\begin{verbatim}
   ... ,"2015","2016","2017","2018",
\end{verbatim}

\hypertarget{check}{%
\subsection{Check}\label{check}}

Before removing it check that all values are \texttt{NA}.

\begin{Shaded}
\begin{Highlighting}[]
\NormalTok{wip }\OperatorTok{%>%}\StringTok{ }\KeywordTok{pull}\NormalTok{(X64) }\OperatorTok{%>%}\StringTok{ }\KeywordTok{is.na}\NormalTok{(.) }\OperatorTok{%>%}\StringTok{ }\KeywordTok{all}\NormalTok{(.)}
\CommentTok{#  [1] TRUE}
\end{Highlighting}
\end{Shaded}

Column \texttt{Indicator.Name} has the unique value \emph{``Proportion
of seats held by women in national parliaments (\%)''} and in
\texttt{Indicator.Code} it is \emph{``SG.GEN.PARL.ZS''}. As there is
only one indicator in this dataset we will remove these two columns.

\hypertarget{exercise-2}{%
\subsection{Exercise}\label{exercise-2}}

Confirm that both \texttt{Indicator.Name} and \texttt{Indicator.Code}
have the same values for all observations. Hint: Use either
\texttt{count()} or \texttt{distinct()} functions.

\hypertarget{removing-columns}{%
\subsection{Removing columns}\label{removing-columns}}

The indicator and \texttt{X64} columns can be removed. We will also
rename ``\texttt{Country.Name}'' as ``\texttt{Country}'' and
``\texttt{Country.Code}'' as ``\texttt{Code}''.

\begin{Shaded}
\begin{Highlighting}[]
\NormalTok{wip2 <-}\StringTok{ }\NormalTok{wip }\OperatorTok{%>%}\StringTok{ }
\StringTok{  }\KeywordTok{select}\NormalTok{(}\OperatorTok{-}\NormalTok{Indicator.Name, }\OperatorTok{-}\NormalTok{Indicator.Code, }
         \OperatorTok{-}\NormalTok{X64) }\OperatorTok{%>%}\StringTok{ }
\StringTok{  }\KeywordTok{rename}\NormalTok{(}\DataTypeTok{Country=}\NormalTok{Country.Name, }\DataTypeTok{Code=}\NormalTok{Country.Code)}
\KeywordTok{head}\NormalTok{(}\KeywordTok{names}\NormalTok{(wip2))}
\CommentTok{#  [1] "Country" "Code"    "X1960"   "X1961"  }
\CommentTok{#  [5] "X1962"   "X1963"}
\KeywordTok{tail}\NormalTok{(}\KeywordTok{names}\NormalTok{(wip2))}
\CommentTok{#  [1] "X2013" "X2014" "X2015" "X2016" "X2017"}
\CommentTok{#  [6] "X2018"}
\end{Highlighting}
\end{Shaded}

\hypertarget{reshape-to-long-format}{%
\subsection{Reshape to long format}\label{reshape-to-long-format}}

We want to transform the data so that for each country the year (column)
data becomes a row. At the same time we will remove the missing data
(with the \texttt{na.rm} option). We will also create a numeric
\texttt{Year} variable and a \texttt{Ratio} of men to women in
parliament\footnote{Uncomment this line if the \texttt{gather()}
  functions gives you character \texttt{pctWiP}.}.

\begin{Shaded}
\begin{Highlighting}[]
\NormalTok{WP <-}\StringTok{ }\NormalTok{wip2 }\OperatorTok{%>%}\StringTok{ }
\StringTok{  }\KeywordTok{gather}\NormalTok{(}\DataTypeTok{key=}\NormalTok{YearC, }\DataTypeTok{value=}\NormalTok{pctWiP,}\KeywordTok{starts_with}\NormalTok{(}\StringTok{"X"}\NormalTok{),}
         \DataTypeTok{na.rm=}\OtherTok{TRUE}\NormalTok{) }\OperatorTok{%>%}\StringTok{ }
\StringTok{  }\KeywordTok{mutate}\NormalTok{(}\DataTypeTok{Year =} \KeywordTok{parse_number}\NormalTok{(YearC),}
         \CommentTok{# pctWiP = as.numeric(pctWiP),}
         \DataTypeTok{Ratio =}\NormalTok{ (}\DecValTok{100}\OperatorTok{-}\NormalTok{pctWiP)}\OperatorTok{/}\NormalTok{pctWiP) }\OperatorTok{%>%}\StringTok{ }
\StringTok{  }\KeywordTok{select}\NormalTok{(Country, Code, Year, pctWiP, Ratio) }\OperatorTok{%>%}\StringTok{ }
\StringTok{  }\KeywordTok{arrange}\NormalTok{(Country, Year)}
\CommentTok{# Look at the contents of WP}
\KeywordTok{glimpse}\NormalTok{(WP)}
\CommentTok{#  Observations: 5,109}
\CommentTok{#  Variables: 5}
\CommentTok{#  $ Country <chr> "Afghanistan", "Afghanis...}
\CommentTok{#  $ Code    <chr> "AFG", "AFG", "AFG", "AF...}
\CommentTok{#  $ Year    <dbl> 1990, 2005, 2006, 2007, ...}
\CommentTok{#  $ pctWiP  <dbl> 3.7, 27.3, 27.3, 27.7, 2...}
\CommentTok{#  $ Ratio   <dbl> 26.03, 2.66, 2.66, 2.61,...}
\end{Highlighting}
\end{Shaded}

\hypertarget{questions}{%
\section{Questions}\label{questions}}

The objective is to look at the geographical and time trends in the
data. We will answer the following questions.

\begin{itemize}
\tightlist
\item
  What are the time trends for Portugal?
\item
  How does Portugal compare to other countries?
\item
  Which countries have the highest percentage of women in parliament by
  year?
\item
  How do continents compare?
\item
  What are the global trends over time?
\end{itemize}

\hypertarget{exercise---without-programming}{%
\subsection{Exercise - Without
Programming}\label{exercise---without-programming}}

\begin{itemize}
\tightlist
\item
  Which country do you think has the highest percentage of women in
  parliament?
\item
  In each continent (i.e.~Africa, Americas, Asia, Europe and Oceania),
  which country has the highest percentage of women in parliament?
\item
  What is the world percentage of women in parliament in 2018?
\end{itemize}

\hypertarget{exploratory-analysis}{%
\section{Exploratory Analysis}\label{exploratory-analysis}}

\hypertarget{select-a-country}{%
\subsection{Select a country}\label{select-a-country}}

This guide explores how Portugal performs over time and compared to
other countries. Before continuing, select another country for yourself
to repeat the examples and do the exercises.

\hypertarget{time-trends-for-portugal}{%
\subsection{Time trends for Portugal}\label{time-trends-for-portugal}}

First look at the raw data.

\begin{Shaded}
\begin{Highlighting}[]
\CommentTok{# Reset tibble print option to see more rows}
\KeywordTok{options}\NormalTok{(}\DataTypeTok{tibble.print_max =} \DecValTok{25}\NormalTok{)}
\NormalTok{WP }\OperatorTok{%>%}\StringTok{ }\KeywordTok{filter}\NormalTok{(Country}\OperatorTok{==}\StringTok{"Portugal"}\NormalTok{)}
\CommentTok{#  # A tibble: 23 x 5}
\CommentTok{#     Country  Code   Year pctWiP Ratio}
\CommentTok{#     <chr>    <chr> <dbl>  <dbl> <dbl>}
\CommentTok{#   1 Portugal PRT    1990    7.6 12.2 }
\CommentTok{#   2 Portugal PRT    1997   13    6.69}
\CommentTok{#   3 Portugal PRT    1998   13    6.69}
\CommentTok{#   4 Portugal PRT    1999   18.7  4.35}
\CommentTok{#   5 Portugal PRT    2000   17.4  4.75}
\CommentTok{#   6 Portugal PRT    2001   18.7  4.35}
\CommentTok{#   7 Portugal PRT    2002   19.1  4.24}
\CommentTok{#   8 Portugal PRT    2003   19.1  4.24}
\CommentTok{#   9 Portugal PRT    2004   19.1  4.24}
\CommentTok{#  10 Portugal PRT    2005   21.3  3.69}
\CommentTok{#  11 Portugal PRT    2006   21.3  3.69}
\CommentTok{#  12 Portugal PRT    2007   28.3  2.53}
\CommentTok{#  13 Portugal PRT    2008   28.3  2.53}
\CommentTok{#  14 Portugal PRT    2009   27.4  2.65}
\CommentTok{#  15 Portugal PRT    2010   27.4  2.65}
\CommentTok{#  16 Portugal PRT    2011   28.7  2.48}
\CommentTok{#  17 Portugal PRT    2012   28.7  2.48}
\CommentTok{#  18 Portugal PRT    2013   28.7  2.48}
\CommentTok{#  19 Portugal PRT    2014   31.3  2.19}
\CommentTok{#  20 Portugal PRT    2015   34.8  1.87}
\CommentTok{#  21 Portugal PRT    2016   34.8  1.87}
\CommentTok{#  22 Portugal PRT    2017   34.8  1.87}
\CommentTok{#  23 Portugal PRT    2018   34.8  1.87}
\end{Highlighting}
\end{Shaded}

\hypertarget{visualisation}{%
\subsubsection{Visualisation}\label{visualisation}}

It is easier to find trends within a plot.

\begin{Shaded}
\begin{Highlighting}[]
\NormalTok{WP }\OperatorTok{%>%}\StringTok{ }
\StringTok{  }\KeywordTok{filter}\NormalTok{(Country}\OperatorTok{==}\StringTok{"Portugal"}\NormalTok{) }\OperatorTok{%>%}\StringTok{ }
\StringTok{  }\KeywordTok{ggplot}\NormalTok{(}\KeywordTok{aes}\NormalTok{(Year, pctWiP)) }\OperatorTok{+}
\StringTok{  }\KeywordTok{geom_line}\NormalTok{() }\OperatorTok{+}\StringTok{ }\KeywordTok{geom_point}\NormalTok{() }\OperatorTok{+}
\StringTok{  }\KeywordTok{scale_y_continuous}\NormalTok{(}\DataTypeTok{limits=}\KeywordTok{c}\NormalTok{(}\DecValTok{0}\NormalTok{, }\DecValTok{50}\NormalTok{)) }\OperatorTok{+}
\StringTok{  }\KeywordTok{ylab}\NormalTok{(}\StringTok{"% Women in Parliament"}\NormalTok{)}
\end{Highlighting}
\end{Shaded}

\begin{center}\includegraphics{WiP-tidyverse_files/figure-latex/PTplot-1} \end{center}

\hypertarget{interpretation}{%
\subsubsection{Interpretation}\label{interpretation}}

In 1990 Portugal had 7.6\% women in parliament (i.e.~12.2 men for each
woman), which increased to 34.8\% (i.e.~1.87 men for each woman) in
2018. This still falls short of 50\% (i.e.~point of gender parity in
parliament).

\hypertarget{exercise-3}{%
\subsection{Exercise}\label{exercise-3}}

For your chosen country look at the time trend data and the plot. What
is your interpretation? How does it compare to Portugal?

\hypertarget{portugal-versus-european-union-eu-countries}{%
\subsection{Portugal versus European Union (EU)
countries}\label{portugal-versus-european-union-eu-countries}}

We selected six EU countries (due to space limitations) for comparison.
It would be better to compare all EU and/or all European countries.

\begin{Shaded}
\begin{Highlighting}[]
\NormalTok{WP }\OperatorTok{%>%}
\StringTok{  }\KeywordTok{filter}\NormalTok{(Country }\OperatorTok{%in%}\StringTok{ }\KeywordTok{c}\NormalTok{(}\StringTok{"Portugal"}\NormalTok{, }\StringTok{"Sweden"}\NormalTok{, }
      \StringTok{"Spain"}\NormalTok{, }\StringTok{"Hungary"}\NormalTok{, }\StringTok{"Romania"}\NormalTok{, }\StringTok{"Finland"}\NormalTok{, }
      \StringTok{"Germany"}\NormalTok{, }\StringTok{"European Union"}\NormalTok{)) }\OperatorTok{%>%}
\StringTok{  }\KeywordTok{ggplot}\NormalTok{(}\KeywordTok{aes}\NormalTok{(Year, pctWiP, }\DataTypeTok{colour=}\NormalTok{Country)) }\OperatorTok{+}
\StringTok{  }\KeywordTok{geom_line}\NormalTok{() }\OperatorTok{+}
\StringTok{  }\KeywordTok{geom_point}\NormalTok{() }\OperatorTok{+}
\StringTok{  }\KeywordTok{scale_x_continuous}\NormalTok{(}\DataTypeTok{breaks=}\KeywordTok{seq}\NormalTok{(}\DecValTok{1990}\NormalTok{, }\DecValTok{2020}\NormalTok{, }\DecValTok{5}\NormalTok{)) }\OperatorTok{+}
\StringTok{  }\KeywordTok{scale_y_continuous}\NormalTok{(}\DataTypeTok{limits=}\KeywordTok{c}\NormalTok{(}\DecValTok{0}\NormalTok{, }\DecValTok{50}\NormalTok{), }
                     \DataTypeTok{breaks=}\KeywordTok{seq}\NormalTok{(}\DecValTok{0}\NormalTok{, }\DecValTok{50}\NormalTok{, }\DataTypeTok{by=}\DecValTok{10}\NormalTok{)) }\OperatorTok{+}
\StringTok{  }\KeywordTok{ggtitle}\NormalTok{(}\StringTok{"Women in Parliament: EU Countries"}\NormalTok{) }\OperatorTok{+}
\StringTok{  }\KeywordTok{ylab}\NormalTok{(}\StringTok{"% Women in Parliament"}\NormalTok{)}
\end{Highlighting}
\end{Shaded}

\begin{center}\includegraphics{WiP-tidyverse_files/figure-latex/euPctPlot-1} \end{center}

\hypertarget{interpretation-1}{%
\subsubsection{Interpretation}\label{interpretation-1}}

Since 2007 Portugal has had more women in parliament than the European
Union average. Hungary and Romania both had a higher percentage of women
in parliament in 1990 (around the end of the Cold War) than they have
had since. The key point to note is that none of these countries reaches
equality between males and females in parliament, although Sweden and
Finland come closest.

\hypertarget{a-couple-of-points-to-note}{%
\subsection{A couple of points to
note}\label{a-couple-of-points-to-note}}

\hypertarget{germany}{%
\subsubsection{``Germany''}\label{germany}}

In October 1990, the process of ``German reunification'' lead to the
creation of Germany, which united the former ``German Democratic
Republic'' (East Germany) and the ``Federal Republic of Germany'' (West
Germany). Therefore, since reunification, the data is presented for the
reunified ``Germany'' only. Careful thought should be given to handling,
analysing and interpreting any pre-reunification data (if available).

\hypertarget{european-union}{%
\subsubsection{``European Union''}\label{european-union}}

The ``European Union'' has changed over time (unlike the ``continent of
Europe''). It started in the 1950s as a block of six European countries
(known as the ``European Community'') and has expanded over the years to
28 countries (with the United Kingdom about to depart). This raises the
question of how the European Union average is calculated. For a given
year, is it calculated based on the actual member states in that year or
on all of the current member states?

\hypertarget{exercises}{%
\subsection{Exercises}\label{exercises}}

Compare the country of your choice to four or five other countries by
plotting a line graph similar to the one above.

\hypertarget{countries-with-the-highest-percentage-of-women-in-parliament}{%
\subsection{Countries with the highest percentage of women in
parliament}\label{countries-with-the-highest-percentage-of-women-in-parliament}}

A quick answer can be obtained by looking at the highest percentages.

\begin{Shaded}
\begin{Highlighting}[]
\NormalTok{WP }\OperatorTok{%>%}\StringTok{ }
\StringTok{  }\KeywordTok{arrange}\NormalTok{(}\OperatorTok{-}\NormalTok{pctWiP) }\OperatorTok{%>%}\StringTok{ }
\StringTok{  }\KeywordTok{head}\NormalTok{(}\DecValTok{10}\NormalTok{)}
\CommentTok{#  # A tibble: 10 x 5}
\CommentTok{#     Country Code   Year pctWiP Ratio}
\CommentTok{#     <chr>   <chr> <dbl>  <dbl> <dbl>}
\CommentTok{#   1 Rwanda  RWA    2013   63.8 0.567}
\CommentTok{#   2 Rwanda  RWA    2014   63.8 0.567}
\CommentTok{#   3 Rwanda  RWA    2015   63.8 0.567}
\CommentTok{#   4 Rwanda  RWA    2016   63.8 0.567}
\CommentTok{#   5 Rwanda  RWA    2017   61.3 0.631}
\CommentTok{#   6 Rwanda  RWA    2018   61.3 0.631}
\CommentTok{#   7 Rwanda  RWA    2008   56.3 0.776}
\CommentTok{#   8 Rwanda  RWA    2009   56.3 0.776}
\CommentTok{#   9 Rwanda  RWA    2010   56.3 0.776}
\CommentTok{#  10 Rwanda  RWA    2011   56.3 0.776}
\end{Highlighting}
\end{Shaded}

\hypertarget{data-speaks}{%
\subsection{Data speaks}\label{data-speaks}}

Are you surprised? Data can be very enlightening.

\hypertarget{highest-percentage-by-year}{%
\subsection{Highest percentage by
year}\label{highest-percentage-by-year}}

Which countries have the highest percentage of women in parliament by
year?

\begin{Shaded}
\begin{Highlighting}[]
\NormalTok{WP }\OperatorTok{%>%}\StringTok{ }
\StringTok{  }\KeywordTok{group_by}\NormalTok{(Year) }\OperatorTok{%>%}\StringTok{ }
\StringTok{  }\KeywordTok{arrange}\NormalTok{(Year, }\OperatorTok{-}\NormalTok{pctWiP) }\OperatorTok{%>%}\StringTok{ }
\StringTok{  }\KeywordTok{filter}\NormalTok{(}\KeywordTok{row_number}\NormalTok{()}\OperatorTok{==}\DecValTok{1}\NormalTok{)}
\CommentTok{#  # A tibble: 23 x 5}
\CommentTok{#  # Groups:   Year [23]}
\CommentTok{#     Country Code   Year pctWiP Ratio}
\CommentTok{#     <chr>   <chr> <dbl>  <dbl> <dbl>}
\CommentTok{#   1 Sweden  SWE    1990   38.4 1.60 }
\CommentTok{#   2 Sweden  SWE    1997   40.4 1.48 }
\CommentTok{#   3 Sweden  SWE    1998   40.4 1.48 }
\CommentTok{#   4 Sweden  SWE    1999   42.7 1.34 }
\CommentTok{#   5 Sweden  SWE    2000   42.7 1.34 }
\CommentTok{#   6 Sweden  SWE    2001   42.7 1.34 }
\CommentTok{#   7 Sweden  SWE    2002   45   1.22 }
\CommentTok{#   8 Rwanda  RWA    2003   48.8 1.05 }
\CommentTok{#   9 Rwanda  RWA    2004   48.8 1.05 }
\CommentTok{#  10 Rwanda  RWA    2005   48.8 1.05 }
\CommentTok{#  11 Rwanda  RWA    2006   48.8 1.05 }
\CommentTok{#  12 Rwanda  RWA    2007   48.8 1.05 }
\CommentTok{#  13 Rwanda  RWA    2008   56.3 0.776}
\CommentTok{#  14 Rwanda  RWA    2009   56.3 0.776}
\CommentTok{#  15 Rwanda  RWA    2010   56.3 0.776}
\CommentTok{#  16 Rwanda  RWA    2011   56.3 0.776}
\CommentTok{#  17 Rwanda  RWA    2012   56.3 0.776}
\CommentTok{#  18 Rwanda  RWA    2013   63.8 0.567}
\CommentTok{#  19 Rwanda  RWA    2014   63.8 0.567}
\CommentTok{#  20 Rwanda  RWA    2015   63.8 0.567}
\CommentTok{#  21 Rwanda  RWA    2016   63.8 0.567}
\CommentTok{#  22 Rwanda  RWA    2017   61.3 0.631}
\CommentTok{#  23 Rwanda  RWA    2018   61.3 0.631}
\end{Highlighting}
\end{Shaded}

\hypertarget{merging-continent}{%
\subsection{Merging continent}\label{merging-continent}}

The variable \texttt{Country} in the \texttt{WP} dataset is a mix of
countries and regions (e.g.~ ``European Union'', ``South Asia'' and
``World''). To present the highest percentages grouped by continent we
need to add it. Luckily, given the large number of R packages available,
we can merge the ``\texttt{continent}'' from the ``\texttt{codelist}''
dataset in the ``\texttt{countrycode}'' package.

\begin{Shaded}
\begin{Highlighting}[]
\CommentTok{# Ensure that 'countrycode' package is installed.}
\CommentTok{# install.packages("countrycode")}
\KeywordTok{library}\NormalTok{(countrycode)}
\NormalTok{cl <-}\StringTok{ }\NormalTok{codelist }\OperatorTok{%>%}\StringTok{ }
\StringTok{  }\KeywordTok{select}\NormalTok{(continent, wb) }\OperatorTok{%>%}\StringTok{ }
\StringTok{  }\KeywordTok{rename}\NormalTok{(}\DataTypeTok{Code =}\NormalTok{ wb, }\DataTypeTok{Continent =}\NormalTok{ continent)}
\NormalTok{cWP <-}\StringTok{ }\NormalTok{WP }\OperatorTok{%>%}\StringTok{ }
\StringTok{  }\KeywordTok{left_join}\NormalTok{(cl, }\DataTypeTok{by =} \StringTok{"Code"}\NormalTok{)}
\end{Highlighting}
\end{Shaded}

\hypertarget{highest-percentages-by-year-and-continent}{%
\subsection{Highest percentages by year and
continent}\label{highest-percentages-by-year-and-continent}}

Which countries have the highest percentages in 1990 and 2018?

\begin{Shaded}
\begin{Highlighting}[]
\NormalTok{cWP }\OperatorTok{%>%}\StringTok{ }
\StringTok{  }\KeywordTok{filter}\NormalTok{(Year }\OperatorTok{%in%}\StringTok{ }\KeywordTok{c}\NormalTok{(}\DecValTok{1990}\NormalTok{, }\DecValTok{2018}\NormalTok{) }\OperatorTok{&}\StringTok{ }
\StringTok{           }\OperatorTok{!}\KeywordTok{is.na}\NormalTok{(Continent)) }\OperatorTok{%>%}\StringTok{ }
\StringTok{  }\KeywordTok{group_by}\NormalTok{(Continent, Year) }\OperatorTok{%>%}\StringTok{ }
\StringTok{  }\KeywordTok{arrange}\NormalTok{(Continent, Year, }\OperatorTok{-}\NormalTok{pctWiP) }\OperatorTok{%>%}\StringTok{ }
\StringTok{  }\KeywordTok{filter}\NormalTok{(}\KeywordTok{row_number}\NormalTok{()}\OperatorTok{==}\DecValTok{1}\NormalTok{) }\OperatorTok{%>%}\StringTok{ }
\StringTok{  }\KeywordTok{select}\NormalTok{(Continent, Country, Year, pctWiP, Ratio)}
\CommentTok{#  # A tibble: 10 x 5}
\CommentTok{#  # Groups:   Continent, Year [10]}
\CommentTok{#     Continent Country        Year pctWiP Ratio}
\CommentTok{#     <chr>     <chr>         <dbl>  <dbl> <dbl>}
\CommentTok{#   1 Africa    Guinea-Bissau  1990   20   4    }
\CommentTok{#   2 Africa    Rwanda         2018   61.3 0.631}
\CommentTok{#   3 Americas  Guyana         1990   36.9 1.71 }
\CommentTok{#   4 Americas  Cuba           2018   53.2 0.880}
\CommentTok{#   5 Asia      Armenia        1990   35.6 1.81 }
\CommentTok{#   6 Asia      Timor-Leste    2018   33.8 1.96 }
\CommentTok{#   7 Europe    Sweden         1990   38.4 1.60 }
\CommentTok{#   8 Europe    Sweden         2018   46.1 1.17 }
\CommentTok{#   9 Oceania   New Zealand    1990   14.4 5.94 }
\CommentTok{#  10 Oceania   New Zealand    2018   38.3 1.61}
\end{Highlighting}
\end{Shaded}

\hypertarget{decline-in-percentage}{%
\subsection{Decline in percentage}\label{decline-in-percentage}}

Which countries have had a decline in percentage since their first
measurement (not always 1990)?

\begin{Shaded}
\begin{Highlighting}[]
\NormalTok{dWP <-}\StringTok{ }\NormalTok{cWP }\OperatorTok{%>%}\StringTok{ }
\StringTok{  }\KeywordTok{group_by}\NormalTok{(Country) }\OperatorTok{%>%}\StringTok{ }
\StringTok{  }\KeywordTok{arrange}\NormalTok{(Country, Year) }\OperatorTok{%>%}
\StringTok{  }\KeywordTok{filter}\NormalTok{(}\KeywordTok{row_number}\NormalTok{()}\OperatorTok{==}\DecValTok{1} \OperatorTok{|}\StringTok{ }\KeywordTok{row_number}\NormalTok{()}\OperatorTok{==}\KeywordTok{n}\NormalTok{()) }\OperatorTok{%>%}\StringTok{ }
\StringTok{  }\KeywordTok{mutate}\NormalTok{(}\DataTypeTok{pctDiff =}\NormalTok{ pctWiP }\OperatorTok{-}\StringTok{ }
\StringTok{           }\KeywordTok{lag}\NormalTok{(pctWiP, }\DataTypeTok{order_by=}\NormalTok{Country)) }\OperatorTok{%>%}\StringTok{ }
\StringTok{  }\KeywordTok{filter}\NormalTok{(pctDiff}\OperatorTok{<}\DecValTok{0}  \OperatorTok{&}\StringTok{ }\OperatorTok{!}\KeywordTok{is.na}\NormalTok{(Continent)) }\OperatorTok{%>%}\StringTok{ }
\StringTok{  }\KeywordTok{arrange}\NormalTok{(pctDiff)}
\NormalTok{dWP }\OperatorTok{%>%}\StringTok{ }\KeywordTok{select}\NormalTok{(Country, pctDiff)}
\CommentTok{#  # A tibble: 17 x 2}
\CommentTok{#  # Groups:   Country [17]}
\CommentTok{#     Country                   pctDiff}
\CommentTok{#     <chr>                       <dbl>}
\CommentTok{#   1 Armenia                   -17.5  }
\CommentTok{#   2 Romania                   -13.7  }
\CommentTok{#   3 Hungary                    -8.1  }
\CommentTok{#   4 Mongolia                   -7.80 }
\CommentTok{#   5 Guinea-Bissau              -6.3  }
\CommentTok{#   6 Guyana                     -5    }
\CommentTok{#   7 Korea, Dem. People’s Rep.  -4.8  }
\CommentTok{#   8 Vanuatu                    -4.3  }
\CommentTok{#   9 Yemen, Rep.                -4.1  }
\CommentTok{#  10 Mali                       -3.40 }
\CommentTok{#  11 Congo, Rep.                -3    }
\CommentTok{#  12 Oman                       -1.2  }
\CommentTok{#  13 Turkmenistan               -1.20 }
\CommentTok{#  14 Haiti                      -1.1  }
\CommentTok{#  15 Tuvalu                     -1    }
\CommentTok{#  16 Albania                    -0.9  }
\CommentTok{#  17 Maldives                   -0.400}
\end{Highlighting}
\end{Shaded}

\hypertarget{visualisation-1}{%
\subsubsection{Visualisation}\label{visualisation-1}}

We will plot the trend lines for countries with at least a 5\% decline.
Note that the ``5\%'' is arbitrarily selected.

\begin{Shaded}
\begin{Highlighting}[]
\CommentTok{# Select the countries to plot}
\NormalTok{dclpct <-}\StringTok{ }\NormalTok{dWP }\OperatorTok{%>%}\StringTok{ }
\StringTok{  }\KeywordTok{filter}\NormalTok{(}\OperatorTok{!}\KeywordTok{is.na}\NormalTok{(Continent) }\OperatorTok{&}\StringTok{ }\NormalTok{pctDiff }\OperatorTok{<=}\StringTok{ }\DecValTok{-5}\NormalTok{) }\OperatorTok{%>%}\StringTok{ }
\StringTok{  }\KeywordTok{pull}\NormalTok{(Country)}

\NormalTok{WP }\OperatorTok{%>%}
\StringTok{  }\KeywordTok{filter}\NormalTok{(Country }\OperatorTok{%in%}\StringTok{ }\NormalTok{dclpct) }\OperatorTok{%>%}\StringTok{ }
\StringTok{  }\KeywordTok{ggplot}\NormalTok{(}\KeywordTok{aes}\NormalTok{(Year, pctWiP, }\DataTypeTok{colour=}\NormalTok{Country)) }\OperatorTok{+}
\StringTok{  }\KeywordTok{geom_line}\NormalTok{() }\OperatorTok{+}
\StringTok{  }\KeywordTok{geom_point}\NormalTok{() }\OperatorTok{+}
\StringTok{  }\KeywordTok{scale_x_continuous}\NormalTok{(}\DataTypeTok{breaks=}\KeywordTok{seq}\NormalTok{(}\DecValTok{1990}\NormalTok{, }\DecValTok{2020}\NormalTok{, }\DecValTok{5}\NormalTok{)) }\OperatorTok{+}
\StringTok{  }\KeywordTok{scale_y_continuous}\NormalTok{(}\DataTypeTok{limits=}\KeywordTok{c}\NormalTok{(}\DecValTok{0}\NormalTok{, }\DecValTok{40}\NormalTok{),}
  \DataTypeTok{breaks=}\KeywordTok{seq}\NormalTok{(}\DecValTok{0}\NormalTok{, }\DecValTok{40}\NormalTok{, }\DataTypeTok{by=}\DecValTok{10}\NormalTok{)) }\OperatorTok{+}
\StringTok{  }\KeywordTok{ggtitle}\NormalTok{(}\StringTok{"Women in Parliament: Decline >=5%"}\NormalTok{) }\OperatorTok{+}
\StringTok{  }\KeywordTok{ylab}\NormalTok{(}\StringTok{"% Women in Parliament"}\NormalTok{)}
\end{Highlighting}
\end{Shaded}

\begin{center}\includegraphics{WiP-tidyverse_files/figure-latex/decline5pct-1} \end{center}

\hypertarget{interpretation-2}{%
\subsubsection{Interpretation}\label{interpretation-2}}

There is a consistent decline between 1990 and 1997 that should be
investigated in collaboration with a subject matter expert to understand
the potential causes.

\hypertarget{ranked-status}{%
\subsection{Ranked status}\label{ranked-status}}

Another way to look at the data is to look at the ranking of countries,
which could be done at a global level or by continent. Nonetheless, the
results should be interpreted with caution and an understanding of the
actual percentages. For example, if most countries were around the 50\%
mark, rankings could be misleading and subject to random fluctuations.

\hypertarget{global-ranks-by-year}{%
\subsection{Global ranks by year}\label{global-ranks-by-year}}

We will rank the countries by year based on the percentage of women in
parliaments. The countries with the highest percentage will be ranked
first and the lowest last. A total for the number of countries with data
is included as it varies by year.

\begin{Shaded}
\begin{Highlighting}[]
\NormalTok{cWPrankG <-}\StringTok{ }\NormalTok{cWP }\OperatorTok{%>%}\StringTok{ }
\StringTok{  }\KeywordTok{filter}\NormalTok{(}\OperatorTok{!}\KeywordTok{is.na}\NormalTok{(Continent)) }\OperatorTok{%>%}\StringTok{ }
\StringTok{  }\KeywordTok{group_by}\NormalTok{(Year) }\OperatorTok{%>%}\StringTok{ }
\StringTok{  }\KeywordTok{mutate}\NormalTok{(}\DataTypeTok{RankG =} \KeywordTok{rank}\NormalTok{(}\OperatorTok{-}\NormalTok{pctWiP), }
         \DataTypeTok{TotalG =} \KeywordTok{n}\NormalTok{()) }
\end{Highlighting}
\end{Shaded}

\hypertarget{global-ranking-portugal}{%
\subsection{Global ranking -- Portugal}\label{global-ranking-portugal}}

\begin{Shaded}
\begin{Highlighting}[]
\NormalTok{cWPrankG }\OperatorTok{%>%}\StringTok{ }
\StringTok{  }\KeywordTok{filter}\NormalTok{(Country}\OperatorTok{==}\StringTok{"Portugal"}\NormalTok{) }\OperatorTok{%>%}\StringTok{ }
\StringTok{  }\KeywordTok{select}\NormalTok{(Country, Year, pctWiP, Ratio, RankG, }
\NormalTok{         TotalG) }\OperatorTok{%>%}\StringTok{ }
\StringTok{  }\KeywordTok{arrange}\NormalTok{(Year)}
\CommentTok{#  # A tibble: 23 x 6}
\CommentTok{#  # Groups:   Year [23]}
\CommentTok{#     Country   Year pctWiP Ratio RankG TotalG}
\CommentTok{#     <chr>    <dbl>  <dbl> <dbl> <dbl>  <int>}
\CommentTok{#   1 Portugal  1990    7.6 12.2   66      138}
\CommentTok{#   2 Portugal  1997   13    6.69  41.5    161}
\CommentTok{#   3 Portugal  1998   13    6.69  49.5    163}
\CommentTok{#   4 Portugal  1999   18.7  4.35  24      154}
\CommentTok{#   5 Portugal  2000   17.4  4.75  34      158}
\CommentTok{#   6 Portugal  2001   18.7  4.35  33      168}
\CommentTok{#   7 Portugal  2002   19.1  4.24  41      161}
\CommentTok{#   8 Portugal  2003   19.1  4.24  46      175}
\CommentTok{#   9 Portugal  2004   19.1  4.24  54      182}
\CommentTok{#  10 Portugal  2005   21.3  3.69  45.5    185}
\CommentTok{#  11 Portugal  2006   21.3  3.69  49.5    189}
\CommentTok{#  12 Portugal  2007   28.3  2.53  28      188}
\CommentTok{#  13 Portugal  2008   28.3  2.53  28      187}
\CommentTok{#  14 Portugal  2009   27.4  2.65  33      187}
\CommentTok{#  15 Portugal  2010   27.4  2.65  34      187}
\CommentTok{#  16 Portugal  2011   28.7  2.48  31      188}
\CommentTok{#  17 Portugal  2012   28.7  2.48  35      188}
\CommentTok{#  18 Portugal  2013   28.7  2.48  39      186}
\CommentTok{#  19 Portugal  2014   31.3  2.19  35      187}
\CommentTok{#  20 Portugal  2015   34.8  1.87  29      188}
\CommentTok{#  21 Portugal  2016   34.8  1.87  27      191}
\CommentTok{#  22 Portugal  2017   34.8  1.87  28      192}
\CommentTok{#  23 Portugal  2018   34.8  1.87  29      192}
\end{Highlighting}
\end{Shaded}

\hypertarget{interpretation-3}{%
\subsubsection{Interpretation}\label{interpretation-3}}

Portugal has generally been ranked in the first quartile (25\%) of
countries in the world, with the fluctuations of its ranking most likely
due to random variation.

\hypertarget{exercise-4}{%
\subsection{Exercise}\label{exercise-4}}

For your chosen country, interpret its ranking over the years. How does
it compare to Portugal?

\hypertarget{continent-ranks-by-year}{%
\subsection{Continent ranks by year}\label{continent-ranks-by-year}}

We will rank the countries by year within a continent based on the
percentage of women in parliaments. The countries with the highest
percentage will be ranked first and the lowest last. A total for the
number of countries with data, within each continent, is included as it
varies by year.

\begin{Shaded}
\begin{Highlighting}[]
\NormalTok{cWPx <-}\StringTok{ }\NormalTok{cWPrankG }\OperatorTok{%>%}\StringTok{ }
\StringTok{  }\KeywordTok{filter}\NormalTok{(}\OperatorTok{!}\KeywordTok{is.na}\NormalTok{(Continent)) }\OperatorTok{%>%}\StringTok{ }
\StringTok{  }\KeywordTok{group_by}\NormalTok{(Continent, Year) }\OperatorTok{%>%}\StringTok{ }
\StringTok{  }\KeywordTok{mutate}\NormalTok{(}\DataTypeTok{RankC =} \KeywordTok{rank}\NormalTok{(}\OperatorTok{-}\NormalTok{pctWiP), }
         \DataTypeTok{TotalC =} \KeywordTok{n}\NormalTok{()) }
\end{Highlighting}
\end{Shaded}

\hypertarget{portugals-ranking-in-europe}{%
\subsection{Portugal's ranking in
Europe}\label{portugals-ranking-in-europe}}

\begin{Shaded}
\begin{Highlighting}[]
\NormalTok{cWPx }\OperatorTok{%>%}\StringTok{ }
\StringTok{  }\KeywordTok{ungroup}\NormalTok{() }\OperatorTok{%>%}\StringTok{ }
\StringTok{  }\KeywordTok{filter}\NormalTok{(Country}\OperatorTok{==}\StringTok{"Portugal"}\NormalTok{) }\OperatorTok{%>%}\StringTok{ }
\StringTok{  }\KeywordTok{select}\NormalTok{(Country, Year, pctWiP, Ratio, RankC, }
\NormalTok{         TotalC) }\OperatorTok{%>%}\StringTok{ }
\StringTok{  }\KeywordTok{arrange}\NormalTok{(Year)}
\CommentTok{#  # A tibble: 23 x 6}
\CommentTok{#     Country   Year pctWiP Ratio RankC TotalC}
\CommentTok{#     <chr>    <dbl>  <dbl> <dbl> <dbl>  <int>}
\CommentTok{#   1 Portugal  1990    7.6 12.2   22       28}
\CommentTok{#   2 Portugal  1997   13    6.69  16.5     39}
\CommentTok{#   3 Portugal  1998   13    6.69  18.5     37}
\CommentTok{#   4 Portugal  1999   18.7  4.35  12       38}
\CommentTok{#   5 Portugal  2000   17.4  4.75  16       38}
\CommentTok{#   6 Portugal  2001   18.7  4.35  15       41}
\CommentTok{#   7 Portugal  2002   19.1  4.24  15       39}
\CommentTok{#   8 Portugal  2003   19.1  4.24  17       41}
\CommentTok{#   9 Portugal  2004   19.1  4.24  21       42}
\CommentTok{#  10 Portugal  2005   21.3  3.69  20       42}
\CommentTok{#  11 Portugal  2006   21.3  3.69  21       44}
\CommentTok{#  12 Portugal  2007   28.3  2.53  15       44}
\CommentTok{#  13 Portugal  2008   28.3  2.53  13       44}
\CommentTok{#  14 Portugal  2009   27.4  2.65  15       44}
\CommentTok{#  15 Portugal  2010   27.4  2.65  15       44}
\CommentTok{#  16 Portugal  2011   28.7  2.48  14       44}
\CommentTok{#  17 Portugal  2012   28.7  2.48  15       43}
\CommentTok{#  18 Portugal  2013   28.7  2.48  17       43}
\CommentTok{#  19 Portugal  2014   31.3  2.19  16       43}
\CommentTok{#  20 Portugal  2015   34.8  1.87  12       43}
\CommentTok{#  21 Portugal  2016   34.8  1.87  11       43}
\CommentTok{#  22 Portugal  2017   34.8  1.87  12       43}
\CommentTok{#  23 Portugal  2018   34.8  1.87  12       43}
\end{Highlighting}
\end{Shaded}

\hypertarget{plot-of-portugals-ranking-in-europe}{%
\subsection{Plot of Portugal's ranking in
Europe}\label{plot-of-portugals-ranking-in-europe}}

Below we reproduce the percentage plot to show how Portugal ranks in
relation to six other European countries. Note that the highest
percentage is ranked first and the lowest last.

\begin{Shaded}
\begin{Highlighting}[]
\NormalTok{cWPx }\OperatorTok{%>%}\StringTok{ }
\StringTok{  }\KeywordTok{filter}\NormalTok{(Country }\OperatorTok{%in%}\StringTok{ }\KeywordTok{c}\NormalTok{(}\StringTok{"Portugal"}\NormalTok{, }\StringTok{"Sweden"}\NormalTok{, }
    \StringTok{"Spain"}\NormalTok{, }\StringTok{"Hungary"}\NormalTok{, }\StringTok{"Romania"}\NormalTok{, }\StringTok{"Finland"}\NormalTok{, }
    \StringTok{"Germany"}\NormalTok{)) }\OperatorTok{%>%}
\StringTok{  }\KeywordTok{ggplot}\NormalTok{(}\KeywordTok{aes}\NormalTok{(Year, RankC, }\DataTypeTok{colour=}\NormalTok{Country)) }\OperatorTok{+}
\StringTok{  }\KeywordTok{geom_line}\NormalTok{() }\OperatorTok{+}
\StringTok{  }\KeywordTok{geom_point}\NormalTok{() }\OperatorTok{+}
\StringTok{  }\KeywordTok{scale_x_continuous}\NormalTok{(}\DataTypeTok{breaks=}\KeywordTok{seq}\NormalTok{(}\DecValTok{1990}\NormalTok{, }\DecValTok{2020}\NormalTok{, }\DecValTok{5}\NormalTok{)) }\OperatorTok{+}
\StringTok{  }\KeywordTok{scale_y_continuous}\NormalTok{(}\DataTypeTok{limits=}\KeywordTok{c}\NormalTok{(}\DecValTok{0}\NormalTok{, }\DecValTok{45}\NormalTok{), }
                     \DataTypeTok{breaks=}\KeywordTok{seq}\NormalTok{(}\DecValTok{0}\NormalTok{, }\DecValTok{45}\NormalTok{, }\DataTypeTok{by=}\DecValTok{10}\NormalTok{)) }\OperatorTok{+}
\StringTok{  }\KeywordTok{ggtitle}\NormalTok{(}\StringTok{"Women in Parliament: Ranked"}\NormalTok{) }\OperatorTok{+}
\StringTok{  }\KeywordTok{ylab}\NormalTok{(}\StringTok{"Rank in Europe"}\NormalTok{)}
\end{Highlighting}
\end{Shaded}

\begin{center}\includegraphics{WiP-tidyverse_files/figure-latex/euRankplot-1} \end{center}

\hypertarget{interpretation-4}{%
\subsubsection{Interpretation}\label{interpretation-4}}

A total of 28 European countries had data in 1990, 39 in 1997 and 43 in
2018. Within Europe, Portugal was typically ranked in the second
quartile (25-50\%) with the fluctuations of its ranking most likely due
to random variation.

\hypertarget{exercise-5}{%
\subsection{Exercise}\label{exercise-5}}

How does your chosen country rank within its continent?

\hypertarget{highest-rank-by-year-and-continent}{%
\subsection{Highest rank by year and
continent}\label{highest-rank-by-year-and-continent}}

Which countries have the highest rank in 1990 and 2018? The answer will
coincide with the highest percentages (see above).

\begin{Shaded}
\begin{Highlighting}[]
\NormalTok{cWPx }\OperatorTok{%>%}\StringTok{ }
\StringTok{  }\KeywordTok{filter}\NormalTok{(Year }\OperatorTok{%in%}\StringTok{ }\KeywordTok{c}\NormalTok{(}\DecValTok{1990}\NormalTok{, }\DecValTok{2018}\NormalTok{) }\OperatorTok{&}\StringTok{ }\NormalTok{RankC}\OperatorTok{==}\DecValTok{1}\NormalTok{) }\OperatorTok{%>%}\StringTok{ }
\StringTok{  }\KeywordTok{arrange}\NormalTok{(Continent, Year) }\OperatorTok{%>%}\StringTok{ }
\StringTok{  }\KeywordTok{select}\NormalTok{(Continent, Year, Country, pctWiP, Ratio) }
\CommentTok{#  # A tibble: 10 x 5}
\CommentTok{#  # Groups:   Continent, Year [10]}
\CommentTok{#     Continent  Year Country       pctWiP Ratio}
\CommentTok{#     <chr>     <dbl> <chr>          <dbl> <dbl>}
\CommentTok{#   1 Africa     1990 Guinea-Bissau   20   4    }
\CommentTok{#   2 Africa     2018 Rwanda          61.3 0.631}
\CommentTok{#   3 Americas   1990 Guyana          36.9 1.71 }
\CommentTok{#   4 Americas   2018 Cuba            53.2 0.880}
\CommentTok{#   5 Asia       1990 Armenia         35.6 1.81 }
\CommentTok{#   6 Asia       2018 Timor-Leste     33.8 1.96 }
\CommentTok{#   7 Europe     1990 Sweden          38.4 1.60 }
\CommentTok{#   8 Europe     2018 Sweden          46.1 1.17 }
\CommentTok{#   9 Oceania    1990 New Zealand     14.4 5.94 }
\CommentTok{#  10 Oceania    2018 New Zealand     38.3 1.61}
\end{Highlighting}
\end{Shaded}

\hypertarget{overall-picture}{%
\subsection{Overall picture}\label{overall-picture}}

What are the trends globally? There are various regions defined in the
World Bank data. We can plot them and highlight the world ``average''.

\begin{Shaded}
\begin{Highlighting}[]
\NormalTok{cWP }\OperatorTok{%>%}\StringTok{ }
\StringTok{  }\KeywordTok{filter}\NormalTok{(}\KeywordTok{is.na}\NormalTok{(Continent)) }\OperatorTok{%>%}
\StringTok{  }\KeywordTok{ggplot}\NormalTok{(}\KeywordTok{aes}\NormalTok{(Year, pctWiP, }\DataTypeTok{group=}\NormalTok{Country)) }\OperatorTok{+}
\StringTok{  }\KeywordTok{geom_line}\NormalTok{() }\OperatorTok{+}
\StringTok{  }\KeywordTok{gghighlight}\NormalTok{(Country}\OperatorTok{==}\StringTok{"World"}\NormalTok{, }
              \DataTypeTok{use_direct_label =} \OtherTok{FALSE}\NormalTok{, }\DataTypeTok{use_group_by =} \OtherTok{FALSE}\NormalTok{) }\OperatorTok{+}
\StringTok{  }\KeywordTok{scale_x_continuous}\NormalTok{(}\DataTypeTok{breaks=}\KeywordTok{seq}\NormalTok{(}\DecValTok{1990}\NormalTok{, }\DecValTok{2020}\NormalTok{, }\DecValTok{5}\NormalTok{)) }\OperatorTok{+}
\StringTok{  }\KeywordTok{scale_y_continuous}\NormalTok{(}\DataTypeTok{limits=}\KeywordTok{c}\NormalTok{(}\DecValTok{0}\NormalTok{, }\DecValTok{40}\NormalTok{), }
                     \DataTypeTok{breaks=}\KeywordTok{seq}\NormalTok{(}\DecValTok{0}\NormalTok{, }\DecValTok{40}\NormalTok{, }\DataTypeTok{by=}\DecValTok{10}\NormalTok{)) }\OperatorTok{+}
\StringTok{  }\KeywordTok{ggtitle}\NormalTok{(}\StringTok{"Women in Parliament: Global Trends"}\NormalTok{) }\OperatorTok{+}
\StringTok{  }\KeywordTok{ylab}\NormalTok{(}\StringTok{"% Women in Parliament"}\NormalTok{)}
\end{Highlighting}
\end{Shaded}

\begin{center}\includegraphics{WiP-tidyverse_files/figure-latex/globalTrends-1} \end{center}

\hypertarget{interpretation-5}{%
\subsubsection{Interpretation}\label{interpretation-5}}

The grey lines show that regardless of how we define region the general
trends are upwards. The ``World'' percentage (black line) increased
between 1997 and 2018. In 2018, women in parliament represented 24\%
(i.e.~a ratio of 3.17 men to each woman), which is still less than half
the level before gender parity can be claimed.

\hypertarget{conclusion}{%
\section{Conclusion}\label{conclusion}}

This guide presented an analysis of the percentage of women in
parliament as a real-life case study for some of the tidyverse package.
Although the format limited what could be presented, we can conclude
that the percentage of women in parliament is increasing but that gender
parity in parliaments is still far-off.

There is a lot more that can be said and discussed about the
limitations, interpretation and potential impact of this data which the
World Bank has nicely
\href{https://databank.worldbank.org/data/reports.aspx?source=2\&type=metadata\&series=SG.GEN.PARL.ZS\#}{summarised}.\footnote{\url{https://databank.worldbank.org/data/reports.aspx?source=2\&type=metadata\&series=SG.GEN.PARL.ZS}.}
You are strongly encouraged to read their discussion for a more complete
understanding.

\begin{center}\includegraphics[width=1\linewidth]{/Users/andreia/Documents/repositories/WiP-tidyverse/images/Women_in_Parliament_rect} \end{center}

%\showmatmethods


\bibliography{pinp}
\bibliographystyle{jss}



\end{document}

